% 原书 p.18–p.21（PDF p.26–p.29）
% >>> p026
\section{哈密顿微分同胚（Hamiltonian Diffeomorphism）}

\begin{keytakeaways}
\begin{itemize}
  \item 几何图像：哈密顿动力学给出系统构造辛同胚的办法——取一个（时间依赖的）函数 $H_t$，
        其哈密顿向量场 $X_{H_t}$ 的流即为哈密顿流，时间 $1$ 映射就是哈密顿微分同胚。
  \item 数学工具：等式 (1.15) $\iota_{X_{H_t}}\omega=-\dd H_t$、流的求导公式 (1.16)、
        Cartan 魔术公式 $\mathcal{L}_X=\dd\iota_X+\iota_X\dd$。
  \item 本节结论：哈密顿微分同胚保辛形式（Liouville 定理 1.7.1），$\Ham(M,\omega)$ 是
        $\Symp(M,\omega)$ 的正规子群（命题 1.7.2），且当 $M$ 闭时还是单群（Banyaga 定理 1.7.1）。
  \item 定量图像：$\Symp_0/\Ham\cong H^1(M;\R)/\Gamma$（Banyaga 定理 1.7.2），
        故通量群 $\Gamma$ 精确度量了「辛同胚」与「哈密顿微分同胚」之间的差距。
\end{itemize}
\end{keytakeaways}

哈密顿动力学提供了系统构造辛同胚的方法。对哈密顿微分同胚的研究是动力系统的一个分支，
它深刻地影响了辛几何的现代发展。

\begin{wdefn}{1.7.1}
对任意光滑函数 $H_t:[0,1]\times M\to\R$（称为\textbf{哈密顿函数}，或简称\textbf{哈密顿量}），
存在唯一的向量场 $X_{H_t}$，称为\textbf{哈密顿向量场}，满足
\begin{equation}
  \iota_{X_{H_t}}\omega=\omega(X_{H_t},\cdot)=-\dd H_t .
  \tag{1.15}
\end{equation}
\end{wdefn}

沿 $X_{H_t}$ 积分，就得到一个流 $\varphi=\{\varphi_H^t\}$，$t\in[0,1]$，称为由 $H$ 生成的
\textbf{哈密顿流}，而其时间 $1$ 映射 $\varphi_H^1$ 称为\textbf{哈密顿微分同胚}。

\begin{wrem}{1.7.1}
\begin{enumerate}[label=(\roman*)]
  \item 按定义，由于 $\varphi_H^t$ 来自向量场 $X_{H_t}$ 的积分，它满足
    \begin{equation}
      \frac{\dd\varphi_H^t}{\dd t}=X_{H_t}\circ\varphi_H^t .
      \tag{1.16}
    \end{equation}
  \item 通过重新标度时间参数 $t$，对任意 $t\in[0,1]$，映射 $\varphi_H^t$ 都是哈密顿微分同胚。
\end{enumerate}
\end{wrem}

下面这个结果说明了为什么哈密顿微分同胚比其他（辛）微分同胚更特殊。

\begin{wprop}{1.7.1}（Liouville）
$(M,\omega)$ 上的任何哈密顿微分同胚 $\varphi$ 都在如下意义下保持辛结构：$\varphi^*\omega=\omega$
（从而 $\varphi\in\Symp(M,\omega)$）。特别地，任何哈密顿微分同胚都保持 $M$ 的体积形式。
\end{wprop}

\begin{wproof}
设 $\varphi=\varphi_H^1$，其中 $H$ 为某个哈密顿函数。我们将证明 $\varphi_t^*\omega=\omega$，
这等价于 $\mathcal{L}_{X_{H_t}}\omega=0$，其中 $X_{H_t}$ 是哈密顿向量场 (1.15)。
由 Cartan 魔术公式，
\[
  \mathcal{L}_{X_{H_t}}\omega
  =\dd\iota_{X_{H_t}}\omega+\iota_{X_{H_t}}\dd\omega
  =\dd(\dd H_t)=0 .
\]
注意我们在本质上用到了 $\omega$ 是闭的这一点（定义 1.1.1 中的条件 (ii)）。
\end{wproof}
% [Tutor Note: 严格按 (1.15) 的符号约定 $\iota_{X_{H_t}}\omega=-\dd H_t$，此处的中间项应为
% $\dd(-\dd H_t)$（仍为 $0$）。原文写作 $\dd(\dd H_t)$，符号上少了一个负号但不影响结论。
% 译文保留原文写法，详见 review_ch01.md。另注：这里还隐含地用了「$\varphi_t^*\omega$ 关于 $t$
% 的导数等于 $\varphi_t^*\mathcal{L}_{X}\omega$」这一标准事实，即流与 Lie 导数的关系
% $\frac{\dd}{\dd t}\varphi_t^*\alpha=\varphi_t^*\mathcal{L}_{X}\alpha$，作者未加说明。]

\begin{wexam}{1.7.1}
对任意开子集 $U\subset(M,\omega)$，若 $H_t|_U=0$，则 $\varphi_H^t|_U=\mathbf{1}$。事实上，
$H_t|_U=0$ 蕴涵 $X_{H_t}|_U=0$。因此 $\varphi_H^t|_U=\varphi_H^0|_U=\mathbf{1}$。作为推论，
在 $(\R^{2n},\omega_{\mathrm{std}})$ 中存在许多非线性的辛同胚，即那些不属于
$\mathrm{Sp}(2n)$ 的辛同胚（因为任何线性辛同胚都整体地改变 $\R^{2n}$）。\hfill$//$
\end{wexam}

% >>> p027
对于一个固定的辛流形 $(M,\omega)$，一旦知道了哈密顿函数 $H_t$，计算其对应的 $X_{H_t}$ 与
$\varphi_H^t$ 通常是直截了当的。

\begin{wexam}{1.7.2}
考虑 $\R^2$ 带 $\omega_{\mathrm{std}}$ 及 $H(x,y)=x$。记
$X_H=A\dfrac{\partial}{\partial x}+B\dfrac{\partial}{\partial y}$，其中 $A,B$ 为某些光滑函数。
则
\[
  \iota_{X_H}\omega_{\mathrm{std}}
  =(\dd x\wedge\dd y)(X_H,\cdot)
  =A\,\dd y-B\,\dd x
  =-\dd H(x,y)=-\dd x .
\]
这说明 $A=0$ 且 $B=1$。因此 $X_H=\dfrac{\partial}{\partial y}$，而 $\varphi_H^1$ 是沿 $y$ 方向的
平移。\hfill$//$
\end{wexam}

\begin{wexer}{1.7.1}
对 $(\R^{2n},\omega_{\mathrm{std}})$，考虑
$H(x_1,y_1,\cdots,x_n,y_n)=\sum\limits_{i=1}^{n}(x_i^2+y_i^2)$。计算所得出的哈密顿流
$\varphi_H^t$。
\end{wexer}

接着，我们研究所有哈密顿微分同胚构成的集合及其群结构。我们将在后面第 6 章 6.1 节看到，
这提供了一个几何视角去探索哈密顿动力系统中的动力学性质，这种视角与源自分析（例如偏微分方程）
的经典方法有根本的不同。

\begin{wdefn}{1.7.2}
如下定义 $(M,\omega)$ 的\textbf{哈密顿群}（Hamiltonian group）：
\[
  \Ham(M,\omega):=\{\varphi\in\Diff(M,\omega)\mid\varphi=\varphi_H^1\}
  \subseteq\Symp(M,\omega).
\]
\end{wdefn}

事实上，$\Ham(M,\omega)\le\Symp_0(M,\omega)$，因为按定义任何哈密顿微分同胚都连回
$\mathbf{1}$。

为使 $\Ham(M,\omega)$ 这一名称名副其实，让我们验证它确实是一个群。为简洁起见，用
$\{\psi_t\},\{\varphi_t\}$ 分别记由 $H_t$ 与 $G_t$ 生成的哈密顿流。于是我们断言
$\{\psi_t\circ\varphi_t\}$ 由
\begin{equation}
  H_t+G_t\circ\psi_t^{-1}
  \tag{1.17}
\end{equation}
生成。

事实上，对任意 $x\in M$，把 $\psi_t\circ\varphi_t$ 关于 $t$ 求导，我们有
\[
  \begin{aligned}
    \frac{\dd}{\dd t}(\psi_t\circ\varphi_t)(x)
    &=\frac{\dd\psi_t}{\dd t}(\varphi_t(x))
      +(\psi_t)_{*,\varphi_t(x)}\left(\frac{\dd\varphi_t}{\dd t}(x)\right)\\
    &=\left(\frac{\dd\psi_t}{\dd t}\circ\psi_t^{-1}\right)(\psi_t\circ\varphi_t)(x)
      +(\psi_t)_{*,\varphi_t(x)}
      \left(\left(\frac{\dd\varphi_t}{\dd t}\circ\varphi_t^{-1}\right)(\varphi_t(x))\right)\\
    &=X_{H_t}\big((\psi_t\circ\varphi_t)(x)\big)
      +(\psi_t)_{*,\varphi_t(x)}\big((X_{G_t})(\varphi_t(x))\big).
  \end{aligned}
\]
这里最后一个等号来自 (1.16)。同时，流 $\{\psi_t\circ\varphi_t\}$ 也给出一个向量场，记作 $X$，
它由下式显式给出
\[
  X\big((\psi_t\circ\varphi_t)(x)\big)=\frac{\dd}{\dd t}(\psi_t\circ\varphi_t)(x).
\]

% >>> p028
因此，通过左乘 $(\psi_t\circ\varphi_t)^{-1}$，我们有
\[
  X(x)=X_{H_t}(x)+(\psi_t)_{*,\psi_t^{-1}(x)}
  \big((X_{G_t})(\psi_t^{-1}(x))\big)\in T_xM .
\]
现在作用 $\omega_x$，我们得到
\[
  \begin{aligned}
    \omega_x(X(x),\cdot)
    &=\omega_x(X_{H_t}(x),\cdot)
      +\omega_x\big((\psi_t)_{*,\psi_t^{-1}(x)}
      ((X_{G_t})(\psi_t^{-1}(x))),\cdot\big)\\
    &=\dd H_t(x)(\cdot)
      +((\psi_t)^*\omega)_{\psi_t^{-1}(x)}
      \big((X_{G_t})(\psi_t^{-1}(x)),(\psi_t^{-1})_*(\cdot)\big)\\
    &=\dd H_t(x)(\cdot)
      +\omega_{\psi_t^{-1}(x)}
      \big((X_{G_t})(\psi_t^{-1}(x)),(\psi_t^{-1})_*(\cdot)\big)\\
    &=\dd H_t(x)(\cdot)
      +\dd G_t(\psi_t^{-1}(x))\big((\psi_t^{-1})_*(\cdot)\big)\\
    &=\dd H_t(x)(\cdot)+\dd(G_t\circ\psi_t^{-1})(x)(\cdot)\\
    &=\dd\big(H_t+G_t\circ\psi_t^{-1}\big)(x)(\cdot).
  \end{aligned}
\]
% [Tutor Note: 原文疑误——按 (1.15) 的约定 $\omega(X_{H_t},\cdot)=-\dd H_t$，这串等式中每一行的
% 「$\dd$」前都应带负号（第一项为 $-\dd H_t$，第二项中 $\omega(Y,\cdot)=-\dd G_t(Y)(\cdot)$），
% 正确的末行应为 $\omega_x(X(x),\cdot)=-\dd(H_t+G_t\circ\psi_t^{-1})(x)(\cdot)$，这才与
% 定义 1.7.1 相符。原文通篇漏写负号，但结论（$X$ 由 $H_t+G_t\circ\psi_t^{-1}$ 生成）不变。
% 译文保留原文写法，详见 review_ch01.md。另外第二步中 $(\psi_t)^*\omega=\omega$ 用到了
% $\psi_t$ 是辛同胚（命题 1.7.1），此处是循环论证的合法使用：命题 1.7.1 已独立证明。]

于是由定义 1.7.1，向量场 $X$ 是由哈密顿函数 $H_t+G_t\circ\psi_t^{-1}$ 生成的哈密顿向量场。
剩下的验证来自下面的练习。

\begin{wexer}{1.7.2}
对哈密顿流 $\{\varphi_H^t\}_{t\in[0,1]}$，证明 $\{\varphi_H^{-t}\}_{t\in[0,1]}$ 也是哈密顿流，
且由 $-H_t\circ\varphi_H^t$ 生成。
\end{wexer}

\begin{wrem}{1.7.2}
\begin{enumerate}[label=(\roman*)]
  \item 关于 (1.17) 有许多直接却富有意义的观察。例如，若 $G_t$ 沿流 $\psi_t$ 不变，
        则 (1.17) 化为 $H_t+G_t$。这特别适用于 $G_t=H_t$ 的情形，因为 $H_t$ 沿其自身的
        哈密顿流不变（由于 $\dd H_t(X_{H_t})=-\omega(X_{H_t},X_{H_t})=0$）。因此，
        任何迭代 $(\varphi_H^t)^k$ 都由 $kH_t$ 生成，$k\in\Z$。
  \item 在练习 1.7.2 中，所得的哈密顿微分同胚是 $\varphi_H^{-1}$。然而，可能存在其他哈密顿流，
        其时间 $1$ 映射同样是 $\varphi_H^{-1}$，但看起来与流 $\{\varphi_H^{-t}\}_{t\in[0,1]}$
        很不一样。例如，考虑 $\widetilde H_t(x):=-H(1-t,x)$。一个标准的计算表明，其哈密顿流为
        $\varphi_{\widetilde H}^t=\varphi_H^{1-t}\circ\varphi_H^{-1}$，时间 $1$ 映射为
        $\varphi_H^{-1}$（见考试 III 中的问题 C (a)）。事实上，流 $\varphi_{\widetilde H}^t$
        可以看作 $\varphi_H^t$ 的时间反演。
\end{enumerate}
\end{wrem}
% [Tutor Note: 原文写作「time-reserving」，应为「time-reversing」（时间反演）之误；译文按其显然
% 的含义译作「时间反演」并在 review_ch01.md 中记录。另注：(i) 中
% $\dd H_t(X_{H_t})=-\omega(X_{H_t},X_{H_t})=0$ 再次用了 (1.15) 的负号约定（斜对称性使
% $\omega(X,X)=0$），故 $H_t$ 沿自身哈密顿流为常数，这正是「能量守恒」的严格表述。]

\begin{wexer}{1.7.3}
设 $H(t,x):[0,1]\times M\to\R$ 是辛流形 $(M,\omega)$ 上生成哈密顿流
$\{\varphi_H^t\}_{t\in[0,1]}$ 的（时间依赖的）函数。证明：对任意时间重参数化
$\alpha(t):[0,1]\to[0,1]$，只要 $\alpha(0)=0$，流 $\{\varphi_H^{\alpha(t)}\}_{t\in[0,1]}$
也是哈密顿流（请写出该流所对应的函数）。并由此推断：若
$\varphi=\varphi_H^1\in\Ham(M,\omega)$，则对任意 $t\in[0,1]$，微分同胚
$\varphi_H^t\in\Ham(M,\omega)$。
\end{wexer}
% [Tutor Note: 原题只要求 $\alpha(0)=0$，但为使 $\{\varphi_H^{\alpha(t)}\}$ 在 $t=1$ 处有意义并
% 得到推论，$\alpha$ 还应满足 $\alpha(1)=1$ 且单调不减（通常取 $\alpha$ 为分段光滑的单调
% 重参数化）。对应函数的写法为 $\alpha'(t)\cdot H(\alpha(t),x)$：由链式法则
% $\frac{\dd}{\dd t}\varphi_H^{\alpha(t)}=\alpha'(t)X_H(\alpha(t),\cdot)$，故生成函数是
% $H$ 沿时间重标度后再乘以 $\alpha'$。由推论立得：$\Ham$ 中元素的任意「时间切片」
% $\varphi_H^t$ 仍是哈密顿微分同胚，这是哈密顿群「路径连通」性质的具体表现。]

\begin{wexer}{1.7.4}
对 $\varphi,\psi\in\Ham(M,\omega)$，若它们在不相交的集合上具有支撑，意即存在哈密顿函数
$H,G:[0,1]\times M\to\R$ 使得 $\varphi_H^1=\varphi$、$\varphi_G^1=\psi$，并且
$\bigcup_{t\in[0,1]}\mathrm{supp}(H_t)$ 与 $\bigcup_{t\in[0,1]}\mathrm{supp}(G_t)$ 不相交，
则 $\varphi\psi=\psi\varphi$。
\end{wexer}

接下来回忆，由命题 1.7.1 有 $\Ham(M,\omega)\le\Symp(M,\omega)$。下面是一个更强的结果。

% >>> p029
\begin{wprop}{1.7.2}
$\Ham(M,\omega)$ 是 $\Symp(M,\omega)$ 的正规子群。具体地，若 $\varphi_t=\varphi_H^t$，则对任意
$\psi\in\Symp(M,\omega)$，共轭 $\psi^{-1}\varphi_t\psi$ 由 $H_t\circ\psi$ 生成。
\end{wprop}

\begin{wproof}
证明与「$\Ham(M,\omega)$ 是群」的验证类似。
\[
  X\circ(\psi^{-1}\varphi_t\psi)
  =\frac{\dd}{\dd t}(\psi^{-1}\varphi_t\psi)
  =(\psi^{-1})_{*,\varphi_t\circ\psi}
  \big(X_{H_t}\circ(\varphi_t\circ\psi)\big),
\]
这蕴涵 $\omega(X,\cdot)=\omega(X_{H_t}\circ\psi,\cdot)=-\dd(H_t\circ\psi)$。因此由定义 1.7.1，
$X$ 是由 $H_t\circ\psi$ 生成的哈密顿向量场。
\end{wproof}

人们可以继续追问：$\Ham(M,\omega)$ 是否含有任何（真）正规子群。辛几何中一个深刻而具有
开创性的定理，由 A. Banyaga [23] 给出，其答案是「否」，只要 $(M,\omega)$ 是闭辛流形。

\begin{wthm}{1.7.1}（Banyaga [23]）
若 $(M,\omega)$ 是闭的，则 $\Ham(M,\omega)$ 是单群。
\end{wthm}

这个定理引出了哈密顿动力学中许多有趣的课题。其中最易入手的一个，是研究 $\Ham(M,\omega)$ 上的
「几乎」同态（称为\textbf{拟同态}，quasi-morphisms）。注意定理 1.7.1 蕴涵 $\Ham(M,\omega)$
上不存在任何非平凡同态，因为否则其核将是 $\Ham(M,\omega)$ 的一个真正规子群。因此，对
$\Ham(M,\omega)$ 做泛函式研究的次优选择，将建立在拟同态之上（见 9.3 节以及 [181] 中的更多
细节）。

沿一个多少相反的方向，人们可以问：$\Ham(M,\omega)$ 与 $\Symp(M,\omega)$（更准确地说，
$\Symp_0(M,\omega)$）之间的差距能有多大？下面这个同样由 Banyaga 给出的定理给出了相当精确的
描述。

\begin{wthm}{1.7.2}（Banyaga [23]）
若 $(M,\omega)$ 是闭的，则
\[
  \Symp_0(M,\omega)/\Ham(M,\omega)\cong H^1(M;\R)/\Gamma,
\]
其中 $\Gamma$ 是 $H^1(M;\R)$ 的一个子群，称为\textbf{通量群}（flux group）。
\end{wthm}

\begin{wexam}{1.7.3}
对 $(S^2,\omega_{\mathrm{area}})$，所有辛同胚都是哈密顿微分同胚，因为
$H^1(S^2;\R)=0$。然而，若 $H^1(M;\R)\neq0$，就可能存在不是哈密顿微分同胚的辛同胚，
如下面的练习所示。\hfill$//$
\end{wexam}

\begin{wexer}{1.7.5}
在标准环面 $(\T^2,\omega_{\mathrm{area}})$ 上，找出一个不是哈密顿微分同胚的辛同胚。并证明
你的答案。
\end{wexer}
